\documentclass[10pt]{article}
\usepackage[a3paper, landscape, margin=1.2cm]{geometry}
\usepackage{longtable}
\usepackage{array}
\usepackage{booktabs}
\usepackage{xcolor}
\usepackage{titlesec}
\usepackage[hidelinks]{hyperref}

\titleformat{\section}{\normalfont\Large\bfseries}{}{0pt}{}
\renewcommand{\arraystretch}{1.25}
\setlength{\tabcolsep}{4pt}
\pagestyle{plain}

\newcolumntype{L}[1]{>{\raggedright\arraybackslash}p{#1}}

\begin{document}

\begin{center}
{\LARGE \textbf{Literature Review: Base Papers on Phishing URL Detection}}\\[4pt]
{\large Extracted directly from each paper's dataset, methodology, results, and limitations sections}
\end{center}
\vspace{6pt}

\footnotesize
\begin{longtable}{L{2.6cm} L{2.6cm} L{1.8cm} L{2.7cm} L{2.7cm} L{4.4cm} L{3.4cm} L{2.6cm} L{3.0cm}}
\toprule
\textbf{Paper} & \textbf{Dataset Used} & \textbf{Size} & \textbf{Phishing URLs} & \textbf{Legitimate URLs} & \textbf{Features Used} & \textbf{Algorithm} & \textbf{Evaluation Metrics} & \textbf{Accuracy} \\
\midrule
\endfirsthead
\toprule
\textbf{Paper} & \textbf{Dataset Used} & \textbf{Size} & \textbf{Phishing URLs} & \textbf{Legitimate URLs} & \textbf{Features Used} & \textbf{Algorithm} & \textbf{Evaluation Metrics} & \textbf{Accuracy} \\
\midrule
\endhead
\bottomrule
\endfoot
\bottomrule
\endlastfoot

ResMLP (Remya et al., IEEE Access 2024)
 & Kaggle dataset (via PhishHaven ref.), filtered to benign+phishing
 & 522{,}214 (filtered; 651{,}191 raw)
 & 94{,}111
 & 428{,}103
 & 25 URL-derived lexical features; WHOIS domain age; Google-index status; GLOVE+NLTK sentiment score
 & ResMLP: conv block + 7 inverted-residual blocks + dense/softmax
 & Accuracy, Precision, Recall, F1, ROC-AUC
 & \textbf{98.29\%} \\

Feature Extension (He et al., IEEE IoT J. 2024)
 & PhishTank (phishing) + Kaggle ``Phishing site URLs'' (legit), balanced
 & 73{,}576
 & 36{,}788
 & 36{,}788
 & TextRank feature-extension library (top-10 keywords) + BERT WordPiece embedding (vector length 138)
 & Custom BERT (6 layers, H=256) $\rightarrow$ 1D-CNN classifier
 & Accuracy, Precision, Recall, F1, FPR, AUC
 & \textbf{99.75\%} \\

Multimodal XAI (Vulfin et al., MAKE 2026)
 & Proprietary D1 + public MTLP (D2)
 & D1 $\approx$3.2M (cleaned); D2 = 15{,}000
 & D1: 1.6M; D2: $\approx$50\% (not printed)
 & D1: 1.6M; D2: $\approx$50\% (not printed)
 & URL+metadata ($>$89 feats incl. WHOIS/SSL); char-level URL (CNN1D); HTML (CodeBERT); screenshot (EfficientNet-B7)
 & Late-fusion ensemble: CatBoost + CNN1D + CodeBERT + EfficientNet-B7 + meta-classifier
 & Accuracy, Precision, Recall, F1, ROC-AUC, FPR, FNR
 & F1 \textbf{0.989} (D1) / \textbf{0.953} (D2) \\

Evaluating Feature Engineering (Kustiawan \& Ghauth, IEEE Access 2025)
 & PhishOFE dataset
 & 101{,}063
 & 48.97\% ($\approx$49{,}491)
 & 51.03\% ($\approx$51{,}572)
 & URL (11 feats) + HTML (19 feats) + Derived (4 feats), tested in 4 combinations
 & 10 ML models compared (RF, DT, kNN, GBM, LR, LightGBM, SVM, CatBoost, NB, XGBoost); best = CatBoost
 & Accuracy, Precision, Recall, F1
 & \textbf{99.45\%} (CatBoost, URL+HTML+derived) \\

PhishOFE (Kustiawan \& Ghauth, IEEE Access 2025)
 & PhishOFE dataset (own; PhishTank+OpenPhish phishing, Open PageRank legit)
 & 101{,}063
 & 48.97\% ($\approx$49{,}491)
 & 51.03\% ($\approx$51{,}572)
 & URL + HTML + derived features; deliberately \textbf{no} third-party/WHOIS features
 & 10 ML models compared; best = CatBoost
 & Accuracy, Precision, Recall, F1, MCC, train/predict time
 & \textbf{99.48\%} (CatBoost) \\

Layered Model (Goenka et al., IEEE Access 2025)
 & Own DS1 (PhishTank+Kaggle) + 6 comparison sets (DS1\_5k/10k/25k, DS2, DS3, DS4)
 & DS1 = 38{,}661
 & 19{,}432
 & 19{,}229
 & Layer 1: 5 brand-jacking features (fuzzy match); Layer 2: 19 lexical features
 & Layer 1: rule-based fuzzy matching; Layer 2: XGBoost (best of 5 ML models)
 & Accuracy, Precision, Recall, F1, response time (ms)
 & \textbf{99.35\%} (12.49 ms); \textbf{89.17\%} on older DS2 (73{,}576 URLs) \\

X-PHIDE (Misiek \& Hyla, IEEE Access 2026)
 & Custom (4 sources) + GramBeddings + PhiUSIIL
 & Custom 75{,}738; GramBeddings 639{,}723 train / 159{,}940 test; PhiUSIIL 235{,}795 (47{,}103 used)
 & PhiUSIIL: 100{,}945; Custom/GramBeddings: $\approx$50\% (rate only)
 & PhiUSIIL: 134{,}850; Custom/GramBeddings: $\approx$50\% (rate only)
 & $>$90 raw features reduced to 8 selected (length\_url, qty\_hyphen\_path, has\_brand\_keyword, subdomain\_auth\_count, qty\_nameservers, qty\_dot\_url, is\_spf\_domain, time\_domain\_activation)
 & XGBoost (Bayesian-tuned)
 & Accuracy, Precision, Recall, Specificity, F1, ROC-AUC, train/inference time, memory
 & \textbf{90.72\%} in-distribution; \textbf{77.84\%} avg.\ cross-dataset \\

Staying Ahead of Phishers (Kavya \& Sumathi, AI Review 2025) --- \textit{survey}
 & N/A --- literature review; cites PhishTank, PhiKitA, UCI repo, Ethereum data across $\sim$50 surveyed studies
 & N/A
 & N/A
 & N/A
 & Surveys list-based, ML, DL (CNN/RNN/LSTM/GRU), GAN, ensemble, NLP, and graph/network-embedding methods
 & N/A (reviews the field, not a single method)
 & Accuracy, Precision, Recall, F1, FPR, AUC-ROC, efficiency (as reported across cited works)
 & N/A --- cited works range \textbf{91.5--98.9\%} \\

Uncovering the Cloak (Li et al., IEEE Access 2023) --- \textit{SLR}
 & N/A --- systematic review of 30 primary studies, 2012--2022
 & N/A
 & N/A
 & N/A
 & Surveys client-side, server-side, URL-manipulation, and content-based cloaking, plus cloaking-detection techniques
 & N/A (reviews the field, not a single method)
 & N/A --- one cited MITM-toolkit classifier reached 99.9\%
 & N/A \\

\end{longtable}

\newpage
\begin{center}
{\Large \textbf{Limitations and Research Gaps, per Paper}}
\end{center}
\vspace{6pt}

\footnotesize
\begin{longtable}{L{4cm} L{10.5cm} L{10.5cm}}
\toprule
\textbf{Paper} & \textbf{Limitation (as stated by the authors)} & \textbf{Gap / Future Work (as stated by the authors)} \\
\midrule
\endfirsthead
\toprule
\textbf{Paper} & \textbf{Limitation (as stated by the authors)} & \textbf{Gap / Future Work (as stated by the authors)} \\
\midrule
\endhead
\bottomrule
\endfoot
\bottomrule
\endlastfoot

ResMLP
 & No dedicated limitations section; implicitly, the 7-block residual architecture adds complexity, and the authors state ``achieving optimal detection accuracy remains an ongoing pursuit.''
 & Extend to multi-class classification; streamline/reduce the residual pipeline's complexity; improve adversarial robustness; integrate real-time monitoring. \\

Feature Extension
 & URL-only: does not analyze page content or malicious JavaScript in phishing sites; BERT component is computationally expensive.
 & Add web-content detection for full coverage; optimize/lighten the BERT model; extend to blockchain and cryptocurrency-specific phishing threats. \\

Multimodal XAI
 & Misclassifies some popular legitimate sites due to coincidental structural similarity; screenshot model sensitive to resolution/aspect-ratio mismatches; late fusion builds no joint cross-modal representation; base models trained independently.
 & Move to a hybrid binding scheme with shared latent representations and cross-modal attention (esp.\ HTML+screenshot); joint cross-modal XAI; systematic adversarial-robustness evaluation. \\

Evaluating Feature Engineering
 & Not explicitly stated for their own study (limitations discussed are of prior/related work, not this paper's).
 & Cascading architecture: cheap URL analysis first, heavier content/relationship analysis only for borderline cases; craft further derived features; ensemble methods with comprehensive feature sets. \\

PhishOFE
 & No dynamic JavaScript/interaction analysis (static features only); not designed for homograph attacks (Latin-script dataset only, not systematically evaluated against them); limited multilingual/internationalized URL coverage; performance tightly coupled to this specific dataset.
 & Homograph-attack detection via character normalization, visual-similarity metrics, and Unicode-aware features; unsupervised learning (clustering, autoencoders, one-class) for novel attacks; multilingual URL datasets; deployable end-user application. \\

Layered Model
 & Brand-jacking layer requires brand names of length $\geq 6$ (shorter names risk false positives); restricted to a fixed list of most-phished brands, so squatting on unlisted brands relies on the lexical layer alone.
 & Expand the brand list; explore alternative transforms for length-based features (binning, rank, log-scale); apply deep learning; extend fuzzy matching to non-space/complex-morphology languages. \\

X-PHIDE
 & Accuracy drops substantially under dataset shift (70.65--90.72\% range); large memory footprint limits browser deployment; no automatic retraining; not built for concurrent-request load; Chrome-only; vulnerable to evolving obfuscation (shorteners, homoglyphs, percent-encoding).
 & Model compression/distillation (pruning, quantization, ONNX); automatic model updating; multi-model ensembling; cross-browser support; adversarial retraining with runtime WHOIS/DNS validation checks. \\

Staying Ahead of Phishers (survey)
 & High false-positive rates in traditional ML/similarity methods; adversarial vulnerability; narrow/imbalanced datasets; poor cross-domain generalizability; ``dark-box'' deep-learning/GAN interpretability; high compute cost; concept drift as attacks evolve.
 & Explainable AI (saliency maps, attention, LIME); federated learning/differential privacy; contextual/behavioral analysis; cross-domain transfer learning; adaptive/incremental retraining; model compression (pruning, quantization, distillation); stronger adversarial defenses. \\

Uncovering the Cloak (SLR)
 & Most detection is server-based, sharing blacklists' latency weakness; very few detection strategies account for cloaking at all; a small number of sources drive over 80\% of successful attacks; crawler ecosystems lack fingerprinting diversity.
 & Broaden temporal scope and databases; diversify literature types reviewed; conduct focused studies of individual cloaking/detection mechanisms; target the small set of sources behind most attacks; improve anti-phishing data-sharing and takedown collaboration. \\

\end{longtable}

\vspace{8pt}
\footnotesize\textit{Note: two papers (Kavya \& Sumathi; Li et al.) are literature reviews rather than
original experiments, and are marked N/A where a column requires a primary dataset, algorithm, or
accuracy figure that a review does not produce on its own.}

\end{document}
